\documentclass[aspectratio=169,xcolor=dvipsnames]{beamer}
\usetheme{Simple}
\usepackage[T1]{fontenc}
\usepackage[utf8]{inputenc}
\usepackage[polish]{babel}
\usepackage{graphicx}
\usepackage{booktabs}
\usepackage{xcolor}
\usepackage{hyperref}

\definecolor{CyberRed}{HTML}{B42332}
\definecolor{CyberBurgundy}{HTML}{7F1D1D}
\definecolor{CyberLight}{HTML}{FCE8EA}
\definecolor{CyberText}{HTML}{292525}

\setbeamercolor{normal text}{fg=CyberText,bg=white}
\setbeamercolor{structure}{fg=CyberRed}
\setbeamercolor{title}{fg=white}
\setbeamercolor{subtitle}{fg=white}
\setbeamercolor{frametitle}{fg=CyberBurgundy,bg=white}
\setbeamercolor{section in toc}{fg=CyberBurgundy}
\setbeamercolor{item}{fg=CyberRed}
\setbeamercolor{subitem}{fg=CyberRed}
\setbeamercolor{enumerate item}{fg=CyberRed}
\setbeamercolor{alerted text}{fg=CyberRed}
\setbeamercolor{block title}{fg=white,bg=CyberBurgundy}
\setbeamercolor{block body}{fg=CyberText,bg=CyberLight}
\setbeamercolor{block title alerted}{fg=white,bg=CyberRed}
\setbeamercolor{block body alerted}{fg=CyberText,bg=CyberLight}
\setbeamercolor{block title example}{fg=white,bg=CyberBurgundy}
\setbeamercolor{block body example}{fg=CyberText,bg=CyberLight}
\setbeamercolor{palette primary}{fg=white,bg=CyberRed}
\setbeamercolor{palette secondary}{fg=white,bg=CyberBurgundy}
\setbeamercolor{palette tertiary}{fg=white,bg=CyberBurgundy}
\setbeamercolor{footline}{fg=CyberBurgundy,bg=white}
\setbeamercolor{navigation symbols}{fg=CyberRed}
\hypersetup{colorlinks=true,urlcolor=CyberRed,linkcolor=CyberBurgundy,citecolor=CyberRed}
\usepackage{tikz}
\usetikzlibrary{positioning,arrows.meta}

% Etykieta prowadzącego slajd (prawy górny róg)
\newcommand{\kto}[1]{%
  \begin{tikzpicture}[remember picture,overlay]
    \node[anchor=north east,fill=CyberLight,text=CyberBurgundy,font=\scriptsize\bfseries,
          rounded corners=2pt,inner sep=3pt] at ([xshift=-3mm,yshift=-3mm]current page.north east) {#1};
  \end{tikzpicture}}

% Ramka na slajdy do uzupełnienia — usunąć przed prezentacją
\newcommand{\dozrobienia}[1]{%
  \begin{alertblock}{Do uzupełnienia}
    \small #1
  \end{alertblock}}

\title[Credential stuffing]{Credential stuffing}
\subtitle{Przejmowanie kont z użyciem danych uwierzytelniających z wycieków\\Moduł 1: Cyberzagrożenia}
\author[Pągowski, Łokaj, Świderski]{Dawid Pągowski \and Bartłomiej Łokaj \and Paweł Świderski}
\institute[PŚ]{Wydział Matematyki Stosowanej\\Politechnika Śląska}
\date{\today}

\begin{document}

\begin{frame}
    \titlepage
\end{frame}

\begin{frame}{Plan prezentacji}
    \tableofcontents
\end{frame}

% =====================================================================
% OSOBA 1 — Dawid Pągowski
% =====================================================================
\section{Wprowadzenie}
\begin{frame}{Cel i zakres analizy}
    \kto{Dawid Pągowski}
    \begin{itemize}
        \item \textbf{Credential stuffing} --- wykorzystywanie wykradzionych poświadczeń
              (np.\ z wycieków) do nieautoryzowanego logowania na konta ofiar.
        \item Kategoria zagrożenia: \textbf{przejęcie konta}.
        \item Zakres:
        \begin{itemize}
            \item definicja i mechanizm, odniesienie do MITRE ATT\&CK
            \item przykład: incydent 23andMe (2023) i skutki w modelu CIA
            \item wykrywanie i zabezpieczenia
        \end{itemize}
        \item Tylko źródła publiczne --- bez testów na cudzych systemach.
    \end{itemize}
\end{frame}

\begin{frame}{Dlaczego ten temat?}
    \kto{Dawid Pągowski}
    \begin{columns}[T]
        \column{0.48\textwidth}
        \centering
        {\Huge\bfseries\color{CyberRed} 19\%}\\[0.1cm]
        \small dziennych prób uwierzytelnienia w logach SSO to credential stuffing
        (mediana)~\cite{verizon}
        \column{0.48\textwidth}
        \centering
        {\Huge\bfseries\color{CyberRed} 49\%}\\[0.1cm]
        \small tyle zapisanych haseł użytkownika jest unikalnych (mediana) ---
        reszta powtarza się w różnych serwisach~\cite{verizon}
    \end{columns}
    \vspace{0.8cm}
    \begin{itemize}
        \item Zagrożenie dobrze udokumentowane przez MITRE i OWASP~\cite{mitre,owasp}.
    \end{itemize}
\end{frame}

\section{Mechanizm}
\begin{frame}{Czym jest credential stuffing?}
    \kto{Dawid Pągowski}
    \begin{block}{Definicja (OWASP)}
        Zautomatyzowane wstrzykiwanie skradzionych poświadczeń (par login--hasło)
        do formularzy logowania w celu uzyskania nieautoryzowanego dostępu do kont~\cite{owasp}.
    \end{block}
    \vspace{0.3cm}
    \begin{itemize}
        \item Podzbiór ataków \textbf{brute force}~\cite{owasp}.
        \item Brute force zgaduje kolejne kombinacje --- credential stuffing próbuje
              \textbf{wyłącznie znanych} poświadczeń.
    \end{itemize}
\end{frame}

\begin{frame}{Zagrożenie, podatność, incydent, ryzyko}
    \kto{Dawid Pągowski}
    \small
    \begin{tabular}{@{}lp{0.78\textwidth}@{}}
        \toprule
        \textbf{Zagrożenie} & zautomatyzowane próby logowania danymi z wycieków na dużą skalę \\
        \textbf{Podatność} & te same hasła w wielu serwisach, brak MFA, brak limitów prób
                             --- słabość procesu i konfiguracji, nie błąd w kodzie \\
        \textbf{Incydent} & udane zalogowanie atakującego na konto ofiary \\
        \textbf{Ryzyko} & wysokie prawdopodobieństwo (tanie listy poświadczeń, automatyzacja)
                          $\times$ wartość danych dostępnych z konta \\
        \bottomrule
    \end{tabular}
\end{frame}

\begin{frame}{Przebieg ataku}
    \kto{Dawid Pągowski}
    \centering
    \begin{tikzpicture}[
        etap/.style={draw=CyberBurgundy,fill=CyberLight,rounded corners=3pt,
                     minimum height=2cm,align=center,
                     font=\footnotesize,text width=3cm,
                     execute at begin node={\hyphenpenalty=10000\exhyphenpenalty=10000}},
        strz/.style={-{Stealth[length=2.5mm]},thick,CyberRed},
        node distance=0.35cm]
        \node[etap] (a) {\textbf{1. Pozyskanie}\\wycieki, phishing, infostealery};
        \node[etap,right=of a] (b) {\textbf{2. Masowe logowanie}\\automatyczne próby na wielu serwisach};
        \node[etap,right=of b] (c) {\textbf{3. Ukrycie ruchu}\\TOR, proxy rezydencjalne};
        \node[etap,right=of c] (d) {\textbf{4. Wykorzystanie}\\przejęte konto, kolejne ofiary};
        \draw[strz] (a) -- (b); \draw[strz] (b) -- (c); \draw[strz] (c) -- (d);
    \end{tikzpicture}
    \vspace{0.5cm}

    \raggedright\small
    \begin{itemize}
        \item Ruch przez proxy rezydencjalne wygląda jak zwykli użytkownicy domowi,
              więc blokady IP niewiele dają~\cite{okta}.
        \item Atakujący może też użyć danych do \textbf{resetu hasła}: odpowiedzi na pytania
              zabezpieczające, numery telefonów, dane osobowe~\cite{owasp}.
    \end{itemize}
\end{frame}

\begin{frame}{Potencjalne cele}
    \kto{Dawid Pągowski}
    \begin{columns}[T]
        \column{0.48\textwidth}
        \begin{block}{Użytkownicy}
            Osoby używające tych samych loginów i haseł w wielu miejscach.
        \end{block}
        \column{0.48\textwidth}
        \begin{block}{Usługi}
            Wszystkie serwisy i aplikacje z publicznym formularzem logowania~\cite{mitre}.
        \end{block}
    \end{columns}
\end{frame}

\begin{frame}{Na tle pokrewnych ataków}
    \kto{Dawid Pągowski}
    \begin{table}
        \centering
        \small
        \begin{tabular}{lll}
            \toprule
            \textbf{Atak} & \textbf{Co jest testowane} & \textbf{Źródło haseł} \\
            \midrule
            Brute force & wiele haseł, jedno konto & słownik / kombinacje \\
            Password spraying & jedno hasło, wiele kont & najpopularniejsze hasła \\
            \textbf{Credential stuffing} & \textbf{konkretne pary login--hasło} & \textbf{wycieki z innych serwisów} \\
            \bottomrule
        \end{tabular}
        \caption{Na podstawie OWASP~\cite{owaspcs}}
    \end{table}
\end{frame}

\begin{frame}{Macierz MITRE ATT\&CK}
    \kto{Dawid Pągowski}
    \begin{table}
        \centering
        \small
        \begin{tabular}{lll}
            \toprule
            \textbf{Taktyka} & \textbf{Technika} & \textbf{Etap} \\
            \midrule
            Reconnaissance & T1589.001 Credentials & 1. Pozyskanie \\
            \textbf{Credential Access} & \textbf{T1110.004 Credential Stuffing} & \textbf{2. Masowe logowanie} \\
            Command and Control & T1090.002 / T1090.003 Proxy & 3. Ukrycie ruchu \\
            Initial Access / Persistence & T1078 Valid Accounts & 4. Wykorzystanie \\
            \bottomrule
        \end{tabular}
        \caption{Na podstawie MITRE ATT\&CK~\cite{mitre}}
    \end{table}
    \small Credential stuffing to zwykle \textbf{jeden etap dłuższego łańcucha}, a nie cel sam w sobie.
\end{frame}

% =====================================================================
% OSOBA 2 — przykład i skutki
% =====================================================================
\section{Przykład i skutki}

\begin{frame}{Przykład: 23andMe (2023)}
    \kto{Paweł Świderski}
    \small
    \begin{itemize}
        \item \textbf{Zasób:} Baza kont użytkowników oraz funkcja \textit{DNA Relatives} (łączenie krewnych).
        \item \textbf{Wejście:} Atak \textit{credential stuffing} (użycie haseł z wycieków w innych serwisach).
        \item \textbf{Działanie:} Przejęcie ok. 14 tys. kont i masowe pobieranie danych połączonych krewnych (głównie osób pochodzenia żydowskiego i chińskiego).
        \item \textbf{Skutek:} Wyciek profili ok. \textbf{6,9 mln} osób (drzewa genealogiczne, pochodzenie, stopień pokrewieństwa).
        \item \textbf{Wykrycie:} Październik 2023 r. – publikacja próbek danych na forach hakerskich.
        \item \textbf{Reakcja:} Wymuszony reset haseł, obowiązkowe uwierzytelnianie 2FA, ugody sądowe.
    \end{itemize}
\end{frame}

\begin{frame}{Skutki w modelu CIA}
    \kto{Paweł Świderski}
    \begin{table}
        \centering
        \footnotesize
        \begin{tabular}{p{2cm} p{2cm} p{6cm}}
            \toprule
            \textbf{Właściwość} & \textbf{Naruszona?} & \textbf{Uzasadnienie} \\
            \midrule
            Poufność & \textbf{Tak} & Ujawnienie danych osobowych i genetycznych 6,9 mln osób. \\
            Integralność & \textbf{Nie} & Dane nie zostały zmodyfikowane w bazie (jedynie skopiowane). \\
            Dostępność & \textbf{Nie} & Serwis działał bez zakłóceń (brak ataku DoS/Ransomware). \\
            \bottomrule
        \end{tabular}
    \end{table}

    \vspace{0.2cm}
    \footnotesize
    \textbf{Pozostałe konsekwencje:}
    \begin{itemize}
        \item \textbf{Finansowe:} Spadek akcji o ponad 90\% (widmo bankructwa), ugoda w USA na 30 mln USD.
        \item \textbf{Prawne:} Naruszenie RODO/UK GDPR, dochodzenia i wielomilionowe kary od regulatorów (np. brytyjskiego ICO).
        \item \textbf{Wizerunkowe:} Trwała utrata zaufania klientów do bezpieczeństwa danych genetycznych.
    \end{itemize}
\end{frame}

% =====================================================================
% OSOBA 3 — wykrywanie, zabezpieczenia, wnioski
% =====================================================================
\section{Wykrywanie i zabezpieczenia}
\begin{frame}{Jak wykryć atak?}
    \kto{Osoba 3}
    \dozrobienia{Sygnały w logach: wiele nieudanych logowań na wiele kont, wiele adresów IP,
    nietypowe lokalizacje, ruch z sieci anonimizujących.}
\end{frame}

\begin{frame}{Proponowane zabezpieczenia}
    \kto{Osoba 3}
    \begin{columns}[T]
        \column{0.32\textwidth}
        \textbf{Zapobieganie}
        \begin{itemize}\small
            \item Uzupełnij
        \end{itemize}
        \column{0.32\textwidth}
        \textbf{Detekcja}
        \begin{itemize}\small
            \item Uzupełnij
        \end{itemize}
        \column{0.32\textwidth}
        \textbf{Reakcja}
        \begin{itemize}\small
            \item Uzupełnij
        \end{itemize}
    \end{columns}
    \vspace{0.3cm}
    \dozrobienia{MFA, sprawdzanie haseł w bazach wycieków, limity prób, CAPTCHA,
    menedżery haseł; reakcja: reset haseł, unieważnienie sesji, powiadomienie użytkowników.}
\end{frame}

\section{Wnioski}
\begin{frame}{Wnioski zespołu}
    \kto{Osoba 3 + zespół}
    \begin{enumerate}
        \item Który etap ataku najłatwiej przerwać?
        \item Które zabezpieczenie daje największy efekt?
        \item Jak sprawdzić, czy ochrona działa?
        \item Jakie są ograniczenia naszej analizy?
    \end{enumerate}
    \dozrobienia{Odpowiedzi na powyższe pytania i rekomendacje zespołu.}
\end{frame}

\section{Źródła}
\begin{frame}{Bibliografia}
    \scriptsize
    \begin{thebibliography}{99}
        \bibitem{owasp} OWASP Foundation, \emph{Credential stuffing},
            \url{https://owasp.org/www-community/attacks/Credential_stuffing}.
        \bibitem{owaspcs} OWASP Foundation, \emph{Credential Stuffing Prevention Cheat Sheet},
            \url{https://cheatsheetseries.owasp.org/cheatsheets/Credential_Stuffing_Prevention_Cheat_Sheet.html}.
        \bibitem{mitre} MITRE, \emph{ATT\&CK: Brute Force: Credential Stuffing (T1110.004)},
            \url{https://attack.mitre.org/techniques/T1110/004/}.
        \bibitem{verizon} P.~Langlois, \emph{Additional 2025 DBIR research on credential stuffing},
            Verizon Business, 2025.
        \bibitem{okta} Help Net Security, \emph{Okta warns customers about credential stuffing
            onslaught}, 29.04.2024.
            \url{https://www.helpnetsecurity.com/2024/04/29/okta-credential-stuffing/}.
        \bibitem{wiki23} Wikipedia, \emph{23andMe data leak}, 
            \url{https://en.wikipedia.org/wiki/23andMe_data_leak}.
        \bibitem{sec23} U.S. Securities and Exchange Commission (SEC), \emph{23andMe Holding Co. Form 8-K}, 
            \url{https://www.sec.gov/Archives/edgar/data/1804591/000119312523253488/d520529d8k.htm}.
        % TODO osoba 2 i 3: dopisać swoje źródła
    \end{thebibliography}
\end{frame}

\begin{frame}{Pytania i dyskusja}
    \centering
    {\LARGE\color{CyberBurgundy} Pytania?}
    \vspace{0.6cm}

    {\large Zapraszamy do dyskusji.}
    % tak naprawde to nie
\end{frame}

\begin{frame}[plain]
    \centering
    {\Huge\bfseries\color{CyberRed} Dziękujemy!}
\end{frame}

\end{document}
